\subsection{STABILIZER, FIXER}

DEFINITION 3. Let M be a monoid operating on a set E and A and B subsets of E. The set of \(\alpha \in \mathbf{M}\) such that \(\alpha \mathbf{A} \subset \mathbf{B}\) (resp. \(\alpha \mathbf{A} = \mathbf{B}\) ) is called the transporter (resp. strict transporter) of A to B. The transporter (resp. strict transporter) of A to A is called the stabilizer (resp. strict stabilizer) of A. The set of \(\alpha \in \mathbf{M}\) such that \(\alpha a = a\) for all \(a \in \mathbf{A}\) is called the fixer of A.

An element \(\alpha\) of M is said to stabilize (resp. stabilize strictly, resp. fix) a subset A of E if \(\alpha\) belongs to the stabilizer (resp. strict stabilizer, resp. fixer) of A. A subset P of M is said to stabilize (resp. stabilize strictly, resp. fix) a subset A of E if all the elements of P stabilize (resp. strictly stabilize, resp. fix) A. The fixer of A is contained in the strict stabilizer of A which itself is contained in the stabilizer of A.

PROPOSITION 1. Let M be a monoid operating on a set E and A a subset of E.

\begin{enumerate}
\def\labelenumi{(\alph{enumi})}
\item
  The stabilizer, strict stabilizer and fixer of A are submonoids of M.
\item
  Let \(\alpha\) be an invertible element of \(\mathbf{M}\) ; if \(\alpha\) belongs to the strict stabilizer (resp. fixer) of \(\mathbf{A}\) , so does \(\alpha^{-1}\) .
\end{enumerate}

Let \(e\) be the identity element of \(\mathbf{M}\) ; then \(ea = a\) for every element \(a \in \mathbf{A}\) and therefore \(e\) belongs to the fixer of \(\mathbf{A}\) . Let \(\alpha\) and \(\beta\) be elements of \(\mathbf{E}\) which stabilize \(\mathbf{A}\) . Then \((\alpha\beta)\mathbf{A} = \alpha(\beta\mathbf{A}) \subset \alpha\mathbf{A} \subset \mathbf{A}\) and therefore the stabilizer of \(\mathbf{A}\) is a submonoid of \(\mathbf{M}\) . Similarly for the strict stabilizer and fixer of \(\mathbf{A}\) , whence (a). If \(\alpha\mathbf{A} = \mathbf{A}\) , then \(\mathbf{A} = \alpha^{-1}(\alpha\mathbf{A}) = \alpha^{-1}\mathbf{A}\) . If for all \(a \in \mathbf{A}\) , \(\alpha a = a\) , then \(a = \alpha^{-1}(\alpha a) = \alpha^{-1}a\) , whence (b).

COROLLARY. Let G be a group operating on a set E and A be a subset of E. The strict stabilizer S and the fixer F of A are subgroups of G and F is a normal subgroup of S.

The first assertion follows from the proposition and F is the kernel of the homomorphism of S into \(S_{A}\) associated with the operation of S on A.

A group G operates faithfully on a set E if and only if the fixer of E consists of the identity element of G. The fixer of E is the kernel of the given homomorphism of G into \(S_{E}\) ; this homomorphism is injective if and only if its kernel consists of the identity element (§ 4, no. 5, Theorem 3).

Let M be a monoid, E an M-set and a an element of E. The fixer, strict stabilizer and stabilizer of \(\{a\}\) are equal; this monoid is called equally the fixer or stabilizer of a. The fixer of a subset A of E is the intersection of the fixers of the elements of A. a is called an invariant element of E if the fixer of a is the monoid M. M is said to operate trivially on E if every element of E is invariant.

PROPOSITION 2. Let G be a group operating on a set E and, for all \(x \in E\) , let \(S_x\) be the stabilizer of \(x\) . For all \(\alpha \in G\) , \(S_{\alpha x} = \alpha S_x \alpha^{-1}\) .

If \(s \in S_x\) , then \(\alpha s \alpha^{-1}(\alpha x) = \alpha s x = \alpha x\) , whence \(\alpha S_x \alpha^{-1} \subset S_{\alpha x}\) . As \(x = \alpha^{-1}(\alpha x)\) , \(\alpha^{-1}S_{\alpha x} \alpha \subset S_x\) , whence \(S_{\alpha x} \subset \alpha S_x \alpha^{-1}\) .

It is seen similarly that, if A and B are two subsets of E and T is the transporter (resp. strict transporter) of A to B, then the transporter (resp. strict transporter) of \(\alpha A\) to \(\alpha B\) is equal to \(\alpha T\alpha^{-1}\) .
% semantic-fragments: legacy-input-seam
\relax{}
